\section{Tests}
\label{sec:tests}

\noindent Pour tester le frontend, en attendant la fin de la conception de l'API, nous avons utilisé une API bouchon.
Celle-ci est implémentée dans le fichier \texttt{apiBack.mock.js}.
Dans ce fichier se trouvent toutes les fonctions de l'API, implémentées de manière à ce qu'elles retournent une réponse "type".
Ainsi, même sans être lié à l'API, nous avons pu coder tous les comportements liés aux réponses de celle-ci.
Cela nous a permis de travailler en parallèle du développement Backend, et de tester différents comportements et cas limites.

\noindent La validation du backend a reposé sur une approche manuelle en deux phases. Les endpoints de l'API REST ont été 
testés via l'interface \textbf{Swagger UI}, générée automatiquement par FastAPI et accessible à l'adresse \texttt{/docs}. 
Grâce à cette interface, nous avons pu tester directement les différentes routes de l'API, en parallèle du développement du 
frontend. Chaque route a été vérifiée dans son cas nominal — création de compte, connexion, calcul d'itinéraire, gestion des 
vélos, des signalements et de l'historique — ainsi que dans quelques cas d'erreur : requête sans token d'authentification 
(\texttt{401 Unauthorized}), corps de requête invalide (\texttt{422 Unprocessable Entity} géré automatiquement par Pydantic), 
et accès à une ressource inexistante (\texttt{404 Not Found}).